\documentclass[10pt,a4paper]{amsart}
\usepackage[margin=19mm]{geometry}
\usepackage{amsmath,amssymb,mathtools}
\usepackage[scr=rsfs,cal=euler]{mathalfa}
\usepackage{xfrac}
\usepackage[unicode,hidelinks,bookmarks=false]{hyperref}
\newcommand{\ie}{i.e.\ }
\newcommand{\cf}{cf.\ }
\newcommand{\resp}{respectively\ }
\newcommand{\Subsection}{\subsection}
\providecommand{\nobreakdash}{\nobreak\hbox{-}}
\setlength{\emergencystretch}{2em}
\multlinegap0pt
\usepackage{mathtools}
\usepackage{amssymb}
\usepackage{xcolor}
\newtheorem{lemma}{Lemma}
\title{Kusner's conjecture is false for $p>4$}
\author{Nathan Xiong}
\date{Source: arXiv:2609.14794v1 (CC BY 4.0)}
\begin{document}
\maketitle

\noindent\emph{Source and licence.} Kusner's conjecture is false for $p>4$. Version arXiv:2609.14794v1 (CC BY 4.0), distributed by arXiv under the Creative Commons Attribution 4.0 International licence, http://creativecommons.org/licenses/by/4.0/, as recorded in the arXiv licence metadata for this exact version and shown on the abstract page https://arxiv.org/abs/2609.14794v1. Full licence notice: \texttt{licenses/arxiv-2609.14794-CC-BY-4.0.txt}.\par\smallskip

\noindent\textit{Editorial scope.} Counted unit: the article's lemma lemma (source0.tex lines 232--235). The article's complete original proof of it is delivered with it (source0.tex lines 237--253, one \texttt{proof} environment of 17 lines). No mathematics is written, repaired or re-derived: the statement and the proof are the article's own lines, in the article's own notation, and no retained line is substituted or re-flowed.  Retained uncounted as the source-local context the proof needs: lines 213--216.  Nothing was omitted from any retained span, so the package carries no omission record.  No retained line contains a citation, so the wrapper carries no bibliography and no bibliography entry.  No retained line contains a figure environment, an image-inclusion command or drawing code, so no image is delivered and the package declares no asset dependency.  Editorial formatting only, in addition: an AMS \texttt{amsart} preamble that restates the article's own declarations and macros used by the retained lines, namely the environments \texttt{lemma} on separate counters, together with the standard packages those lines need.  The retained environments print in this document as Lemma 1 in order of appearance, whereas the article prints the same items with the numbering of its own full document.  The complete native source of the article as distributed in the arXiv e-print for this exact version is bundled unchanged as \texttt{sources/210358/source0.tex}.
\begin{lemma}\label{lem:power-ineq}
For $p>4$ and distinct $x,y>0$,
\[(x+y)^p+|x-y|^p>2x^p+2y^p+(2^p-4)(xy)^{p/2}.\]
\end{lemma}

\begin{lemma}
Let $p>4$. For any sufficiently large $m$, there exists $a>1$ such that, with $R\coloneqq a^p+2$, $A\coloneq (a+1)^p+(a-1)^p+2$, and $B\coloneq 2a^p+2^p$, we have
\[2A-B=2^{p-1}R\left(1+\frac{1}{m}\right).\]
\end{lemma}

\begin{proof}
Set $a_0=4^{1/p}>1$. Applying Lemma~\ref{lem:power-ineq} with
$x=a_0$ and $y=1$ gives
\[(a_0+1)^p+(a_0-1)^p>2^{p+1}+2.\]
Thus, at $a=a_0$, we have $R=6$ and
\begin{align*}
    2A-B &= 2(a_0+1)^p+2(a_0-1)^p+4-2a_0^p-2^p\\
    &> 3\cdot 2^p\\
    &= 2^{p-1}R
\end{align*}
Define the continuous function $\Phi$ over $a\in[1,a_0]$ by
\[\Phi(a)=\frac{2A-B}{2^{p-1}R}.\]
We already showed that $\Phi(a_0)>1$. Then, we can plug in $a=1$ to get
\[\Phi(1)=\frac{2^p+2}{3\cdot2^{p-1}}<1.\]
Therefore, as long as $m>\frac{1}{\Phi(a_0)-1}$, the intermediate value theorem supplies an $a\in(1,a_0)$ such that
$\Phi(a)=1+\frac{1}{m}$.
\end{proof}

\end{document}
